\documentclass[10pt,a4paper]{amsart}
\usepackage[margin=19mm]{geometry}
\usepackage{amsmath,amssymb,mathtools}
\usepackage[scr=rsfs,cal=euler]{mathalfa}
\usepackage{xfrac}
\usepackage[unicode,hidelinks,bookmarks=false]{hyperref}
\newcommand{\ie}{i.e.\ }
\newcommand{\cf}{cf.\ }
\newcommand{\resp}{respectively\ }
\newcommand{\Subsection}{\subsection}
\providecommand{\nobreakdash}{\nobreak\hbox{-}}
\setlength{\emergencystretch}{2em}
\multlinegap0pt
\usepackage{amssymb}
\usepackage{mathtools}
\usepackage{xcolor}
\usepackage[normalem]{ulem}
\usepackage[shortlabels]{enumitem}
\newtheorem{lemma}{Lemma}
\newtheorem{proposition}{Proposition}
\def\Im{\mathop\mathrm{Im}}
\newcommand{\R}{\mathbb{R}}
\def\cal{\mathcal}
\def\cd{\centerdot}
\def\d{\partial}
\newcommand*{\defeq}{\mathrel{\vcenter{\baselineskip0.5ex \lineskiplimit0pt
			\hbox{\scriptsize.}\hbox{\scriptsize.}}}%
	=}
\newcommand{\dist}{\mathrm{dist}}
\def\lec{\lesssim}
\def\ovl{\overline}
\def\wt{\widetilde}
\title{Stability estimates for the initial-to-final-state inverse problem}
\author{Manuel Ca\~nizares, Thanasis Zacharopoulos}
\date{Source: arXiv:2609.00257v1 (CC BY 4.0)}
\begin{document}
\maketitle

\noindent\emph{Source and licence.} Stability estimates for the initial-to-final-state inverse problem. Version arXiv:2609.00257v1 (CC BY 4.0), distributed by arXiv under the Creative Commons Attribution 4.0 International licence, http://creativecommons.org/licenses/by/4.0/, as recorded in the arXiv licence metadata for this exact version and shown on the abstract page https://arxiv.org/abs/2609.00257v1. Full licence notice: \texttt{licenses/arxiv-2609.00257-CC-BY-4.0.txt}.\par\smallskip

\noindent\textit{Editorial scope.} Counted unit: the article's proposition distance (source0.tex lines 789--803). The article's complete original proof of it is delivered with it (source0.tex lines 806--886, one \texttt{proof} environment of 81 lines). No mathematics is written, repaired or re-derived: the statement and the proof are the article's own lines, in the article's own notation, and no retained line is substituted or re-flowed.  Retained uncounted as the source-local context the proof needs: lines 630--635; lines 677--679; lines 696--727; lines 716--718; lines 722--724.  Nothing was omitted from any retained span, so the package carries no omission record.  No retained line contains a citation, so the wrapper carries no bibliography and no bibliography entry.  No retained line contains a figure environment, an image-inclusion command or drawing code, so no image is delivered and the package declares no asset dependency.  Editorial formatting only, in addition: an AMS \texttt{amsart} preamble that restates the article's own declarations and macros used by the retained lines, namely the environments \texttt{lemma}, \texttt{proposition} on separate counters, together with the standard packages those lines need.  The retained environments print in this document as Lemma 1, Proposition 1 in order of appearance, whereas the article prints the same items with the numbering of its own full document.  The complete native source of the article as distributed in the arXiv e-print for this exact version is bundled unchanged as \texttt{sources/209010/source0.tex}.
\begin{equation}\label{eq: distance}
\dist_w(\cal{U}_T^1, \cal{U}_T^2) \defeq
\sup_{\substack{u_1 \in \cal{P}_{V_1}\\ u_2 \in \cal{P}_{\overline{V}_2}}} \frac{\Big| \displaystyle \int_\Sigma (\cal{U}_T^1-\cal{U}_T^2) u_1(0, \cd) \ovl{u_2(T, \cd)}\Big|}{\|u_1\|_{L^2_w(\Sigma)} \| u_2\|_{L^2_w(\Sigma)}}
%
=\sup_{\substack{u_1 \in \cal{P}_{V_1}\\ u_2 \in \cal{P}_{\overline{V}_2}}} \frac{\Big| \displaystyle \int_\Sigma (V_1 - V_2) u_1 \ovl{u_2} \Big|}{\|u_1\|_{L^2_w(\Sigma)} \| u_2\|_{L^2_w(\Sigma)}},
\end{equation}

\begin{equation}\label{eq: bound_IF_map_distance}
	\|\cal{U}_T^1 - \cal{U}^2_T\|_{\cal{L}(L^2(\R^n))} \lec \dist_{w}(\cal{U}_T^1, \cal{U}_T^2). 
\end{equation}

\begin{lemma}\label{lem: closeness_of_solutions_time}
Let $V_j\in L^1((0, T); L^\infty(\R^n))$ with $n\geq 2$ be such that $\| V_j \|_{w}< \infty$ with weight $w:\R^n \to (0,1]$ and satisfying the smallness condition $\|  \Im V_j \|_{w} \leq \frac{1}{4T} $, for $j\in\{1,2,3\}$.
For $f,g\in L^2(\R^n)$, let $u_j, \, v_j \in C([0, T]; L^2(\R^n))$ be solutions of the corresponding initial and final value problems for the Schr\"odinger equation
%
\[
\begin{dcases}
(i\d_t+\Delta - V_j)u_j = 0, \quad &\textup{in}\,\, \Sigma, \\
u_j(0, \cd) = f \in L^2(\R^n), &\textup{in}\,\, \R^n
\end{dcases} 
%
\quad \textup{and} \quad
%
\begin{dcases}
(i\d_t+\Delta - \ovl{V_j})v_j = 0, \quad &\textup{in}\,\, \Sigma, \\
v_j(T, \cd) = g\in L^2(\R^n), &\textup{in}\,\, \R^n.
\end{dcases} 
\]
%
Then, it holds that 
%
\begin{equation}\label{eq: solutions_diff1_time}
\|u_j - u_k \|_{L^2_w(\Sigma)} \leq 4T \| V_j - V_k \|_{w} \|u_k\|_{L^2_w(\Sigma)}
\end{equation}
%
and
%
\begin{equation}\label{eq: solutions_diff2_time}
\|v_j - v_k \|_{L^2_w(\Sigma)} \leq 4T \| V_j - V_k \|_{w} \|v_k\|_{L^2_w(\Sigma)},
\end{equation}
%
for every $j \neq k \in \{1, 2, 3 \}$.
\end{lemma}

\begin{equation}
\|u_j - u_k \|_{L^2_w(\Sigma)} \leq 4T \| V_j - V_k \|_{w} \|u_k\|_{L^2_w(\Sigma)}
\end{equation}

\begin{equation}
\|v_j - v_k \|_{L^2_w(\Sigma)} \leq 4T \| V_j - V_k \|_{w} \|v_k\|_{L^2_w(\Sigma)},
\end{equation}

\begin{proposition}\label{prop: distance}
Let $n\geq 2$ and $V_1, \, V_2 \in L^1((0, T); L^\infty(\R^n))$ with $\| V_j \|_w < \infty$, for $j\in\{1, 2\}$ and weight $w:\R^n \to (0, 1)$. The quantity defined in \eqref{eq: distance} is a distance. 

In particular, fix $R_0>0$ and for $V \in L^1((0, T); L^\infty(\R^n))$ with $\| V \|_w < \infty$ and $\| \Im {V} \|_w\leq \frac{1}{4T}$ set 
%
\[
\cal {T}_V := \{\wt{ V} \in  L^1((0, T); L^\infty(\R^n)) : \,\, \| \wt{V} \|_w<\infty, \, \| \Im \wt{V} \|_w \leq \frac{1}{4T} \,\,
%
\text{and such that} \,\, 
%
\| V - \wt{V} \|_w \leq \frac{R_0}{2} \}.
\]
%
Then, if $\cal J$ is the set of the initial-to-final-state maps corresponding to potentials in $ \cal T_V$, the quantity \eqref{eq: distance} becomes a quasi-metric on the set $\cal J$.
\end{proposition}

\begin{proof}
Obviously, $\dist_{w}(\cal{U}_T^1, \cal{U}_T^2) \geq 0$ and $\dist_{w}(\cal{U}_T^1, \cal{U}_T^2) = \dist_{w}(\cal{U}_T^2, \cal{U}_T^1)$. In order to see that 
%
\[
\dist_{w}(\cal{U}_T^1, \cal{U}_T^2) = 0 \Longleftrightarrow \cal{U}_T^1 = \cal{U}_T^2,
\]
%
first notice that if $\dist_{w}(\cal{U}_T^1, \cal{U}_T^2) = 0$, then by \eqref{eq: bound_IF_map_distance} we have that $\| \cal{U}_T^1 - \cal{U}_T^2\|_{\cal{L}(L^2(\R^n))} = 0$ which readily implies $\cal{U}_T^1=\cal{U}_T^2$. 
%
In the case that $\cal{U}_T^1 = \cal{U}_T^2$, then by definition we immediately have that $\dist_{w}(\cal{U}_T^1, \cal{U}_T^2) = 0$. 
Hence, \eqref{eq: distance} is a distance.
%%

%%%%%
It remains to prove that $\dist_w(\cdot, \cdot)$ satisfies a quasi-triangle inequality under the a priori assumption for the potentials. 
To that end, consider the following initial and final value problems for the Schr\"odinger equation:
%
\[
\begin{dcases}
(i\d_t+\Delta - V_j)u_j = 0, \quad &\textup{in}\,\, \Sigma, \\
u_j(0, \cd) = f \in L^2(\R^n), &\textup{in}\,\, \R^n
\end{dcases} 
%
\quad \textup{and} \quad
%
\begin{dcases}
(i\d_t+\Delta - \ovl{V_j})v_j = 0, \quad &\textup{in}\,\, \Sigma, \\
v_j(T, \cd) = g\in L^2(\R^n), &\textup{in}\,\, \R^n,
\end{dcases} 
\]
%
for $ j\in \{1, 2, 3\}$. Of course $V_j \in L^1((0, T); L^\infty(\R^n))$ for $j\in\{ 1, 2, 3 \}$. 
%
It holds that 
%
\[
i \int_{\R^n} (\cal{U}^1_T -\cal{U}^2_T) f \ovl{g} = i\int_{\R^n} (\cal{U}^1_T - \cal{U}^3_T)f\ovl{g} + i \int_{\R^n}(\cal{U}_T^3-\cal{U}^2_T)f \ovl{g}.
\]
%
Assume also that $\|V_3\|_w<\infty$, $\| \Im V_3 \|_w \leq \frac{1}{4T}$ and $V_1, V_2 \in \cal{T}_{V_3}$.
By taking absolute values and dividing the above equality with the quantity $\|u_1\|_{L^2_w(\Sigma)} \|v_2\|_{L^2_w(\Sigma)}$, we infer that
%
\begin{align*}
&\frac{\Big| \displaystyle \int_{\R^n} (\cal{U}^1_T -\cal{U}^2_T) u_1(0, \cd) \ovl{v_2(T, \cd)}\Big|}{\|u_1\|_{L^2_w(\Sigma)} \|v_2\|_{L^2_w(\Sigma)}} \\
%
&\leq \frac{\Big| \displaystyle \int_{\R^n} (\cal{U}^1_T -\cal{U}^3_T) u_1(0, \cd) \ovl{v_3(T, \cd)}\Big|}{\|u_1\|_{L^2_w(\Sigma)} \|v_3\|_{L^2_w(\Sigma)}} \frac{\|v_3\|_{L^2_w(\Sigma)}}{\|v_2\|_{L^2_w(\Sigma)}}
%
+ \frac{\Big| \displaystyle \int_{\R^n} (\cal{U}^3_T -\cal{U}^2_T) u_3(0, \cd) \ovl{v_2(T, \cd)}\Big|}{\|u_3\|_{L^2_w(\Sigma)} \|v_2\|_{L^2_w(\Sigma)}} \frac{\|u_3\|_{L^2_w(\Sigma)}}{\|u_1\|_{L^2_w(\Sigma)}}
\end{align*}
%
Using the estimate \eqref{eq: solutions_diff1_time} of Lemma \ref{lem: closeness_of_solutions_time} and the fact that $V_1\in\cal{T}_{V_3}$, we have that 
%
\[
\frac{\|u_3\|_{L^2_w(\Sigma)}}{\|u_1\|_{L^2_w(\Sigma)}} \leq \frac{\|u_3 - u_1\|_{L^2_w(\Sigma)}}{\|u_1\|_{L^2_w(\Sigma)}} + 1 
\leq 4T \| V_3 - V_1\|_w +1 
\leq  2T R_0 +1 .
\]
%
Similarly, using \eqref{eq: solutions_diff2_time} instead and that $V_2\in \cal{T}_{V_3}$, we also get 
%
\[
\frac{\|v_3\|_{L^2_w(\Sigma)}}{\|v_2\|_{L^2_w(\Sigma)}} \leq 4T \| V_3 - V_2 \|_w +1
\leq 2T R_0 +1.
\]
%
Setting $ M:=  2T R_0 +1 < \infty$, we obtain 
%
\[
\frac{\Big| \displaystyle \int_{\R^n} (\cal{U}^1_T -\cal{U}^2_T) u_1(0, \cd) \ovl{v_2(T, \cd)}\Big|}{\|u_1\|_{L^2_w(\Sigma)} \|v_2\|_{L^2_w(\Sigma)}} 
\leq M \Big( \dist_w(\cal{U}_T^1, \cal{U}_T^3) + \dist_w(\cal{U}_T^3, \cal{U}_T^2) \Big).
\]
%
By taking supremum over all solutions $u_1\in \cal{P}_{V_1}$, $v_2 \in \cal{P}_{\ovl{V_2}}$, this last estimate implies 
%
\[
\dist_w(\cal{U}_T^1, \cal{U}_T^2) \leq M \Big( \dist_w(\cal{U}_T^1, \cal{U}_T^3) + \dist_w(\cal{U}_T^3, \cal{U}_T^2) \Big)
\]
%
with $M > 1$, which concludes the proof that $\dist_w(\cdot, \cdot)$ is a quasi-metric on the set $\cal{J} \subset \cal{L}(L^2(\R^n))$ of the initial-to-state-maps that correspond to the potentials in $\cal{T}_{V_3}$. 
%
\end{proof}

\end{document}
